% Drafted 2026-08-17. NOT yet inputted by paper_journal/main.tex.
% WHEN WIRING IN: the sparsity subsection of main.tex carries no \label. Two
% places here refer to it in prose instead. Add \label{sec:sparsity} there and
% convert them to \ref{} if cross-references are wanted.
% Numbers trace to experiments/results/annotation_kg/ and to the RESEARCH_LOG.md
% entry dated 2026-08-17 (Step 4). Detector frozen in commit b282957 before any
% held-out answer was read.

\section{Does the Knowledge Graph Change What the Model Recommends?}
\label{sec:kg-violation}

Section~\ref{sec:results} shows that knowledge-graph augmentation buys quality
at substantial latency cost, and our sparsity analysis shows its benefit
shrinking rather than growing with record sparsity. Neither asks the question a
clinician would ask first: when we inject guideline facts into the prompt, does
the model's recommendation actually respect them? We report that check here.

\subsection{Only one relation supports a well-posed check}

Our graph carries five relation types, and they are not equally checkable. Three
of them---first-line treatments, characteristic symptoms, and indicated
laboratory tests---are retrieved as truncated lists of at most three items.
Naming something outside such a list contradicts nothing, because the list was
never exhaustive, and omission is equally uninformative. A fourth,
laboratory co-occurrence, is a frequency computed over admissions; it is a
descriptive statistic, and no clinical assertion can contradict it.

Contraindication is different in kind. It is a prohibition, so a single
endorsing mention is a definite violation, with no exhaustiveness assumption
required. It is also the only relation our retriever fetches without a limit,
and the only one with direct safety consequences. We therefore scope the
measure to it and name it for what it measures: the
\textbf{contraindication-violation rate}, not a general graph-faithfulness
score.

We count a violation when an answer endorses use of a drug the graph marks
contraindicated for \emph{any} condition matched to that patient, not only the
condition named in the question---a drug may be appropriate for the disease
asked about and prohibited by a comorbidity. Restating the graph's own warning
is compliance, not violation. Crucially, the graph facts serve as ground truth
for \emph{all three} retrieval modes, including the two that never receive
them, which is what makes the comparison below a test of whether injecting the
facts changes behaviour.

\subsection{Detector, frozen before use}

Detection reduces to two steps: does the answer name a prohibited drug, and
does it endorse or warn against it. The first is a closed-vocabulary lookup
over fourteen drug names taken from the graph, which avoids the open-vocabulary
matching that made our earlier contradiction heuristic unusable
(Section~\ref{sec:halluc-validation}). The second is polarity, which is where
lexical methods fail.

We wrote the polarity rules against a development set drawn from the
router-training and router-validation splits---data disjoint from the held-out
evaluation---and froze them before reading a single held-out answer. That
discipline paid immediately. The development set contained the answer
\emph{``TSH, Lithium Levels.''}, which names a drug prohibited for that patient
while recommending nothing at all: the drug name is part of a laboratory test
name. This is the same class of error as the coded-vocabulary problem that
invalidated our previous detector, in a form we would not have anticipated from
the schema, and it is now an explicit exclusion. The detector emits four
verdicts---violation, compliant, not applicable, and abstain---and attaches the
matched drug, the graph edge, and the deciding span to every decision so that
each can be audited without re-deriving it.

\subsection{Validation}

A human annotator, blind to the detector's output, labelled 61 generations
stratified over the detector's own verdicts, together with constructed
violations as attention checks. Of 41 generations the detector flagged,
\textbf{41 were confirmed genuine and none were false positives}: precision
$1.000$, exact 95\% CI $[0.914, 1.000]$, against a bar of $0.80$ registered
before annotation. Ten flagged as compliant were confirmed compliant. Of ten
where a prohibition applied but the detector stayed silent, nine were confirmed
silent correctly. The detector abstained on none of the 900 generations.

The single disagreement is instructive rather than a defect. The annotator
judged an endorsement of empagliflozin to be a violation on the basis of that
drug's real renal risk profile; the graph does not encode that contraindication
for the patient concerned, so the detector, applying its stated definition,
returned not applicable. We record this as a \emph{gap in the knowledge graph}
rather than a detector error, and retain recall at $1.000$.

\subsection{Result: the graph repairs damage the record introduces}

Table~\ref{tab:kg-violation} reports violations by retrieval mode over all 300
held-out questions, of which 212 engage at least one prohibition. Comparisons
are paired McNemar tests on the same questions.

\begin{table}[t]
\centering
\footnotesize
\setlength{\tabcolsep}{5pt}
\caption{Contraindication violations by retrieval mode (300 held-out questions), with paired McNemar tests. Adding the record snapshot increases violations; adding graph facts removes them.}
\label{tab:kg-violation}
\begin{tabular}{lc}
\toprule
Mode & Violations \\
\midrule
T (text only)            & 12 \\
T+E (text + record)      & \textbf{18} \\
T+E+K (text + record + graph) & \textbf{11} \\
\midrule
Comparison & Discordant, $p$ \\
\midrule
T+E vs.\ T+E+K & 7 vs.\ 0, $p = 0.016$ \\
T vs.\ T+E     & 0 vs.\ 6, $p = 0.031$ \\
T vs.\ T+E+K   & 1 vs.\ 0, $p = 1.000$ \\
\bottomrule
\end{tabular}
\end{table}

Injecting graph facts eliminates seven violations and introduces none
($p = 0.016$). Read alone, that is a clean safety benefit. Read together with
the other two comparisons, it is something more specific and more interesting.

Adding the record snapshot to text-only retrieval \emph{creates} six violations
and removes none ($p = 0.031$). The mechanism is visible in the failure cases:
presented with a structured medication list, the model affirms a drug it can
see the patient is taking, without reference to whether that drug is
contraindicated by another of the patient's conditions. The record supplies
what the patient \emph{is} receiving, not what they \emph{should} receive, and
the model conflates the two.

Against text-only retrieval, the full hybrid configuration is not significantly
safer ($p = 1.000$). \textbf{The graph's contribution here is therefore
corrective rather than additive: it repairs a specific failure mode introduced
by the record snapshot, rather than improving on the text-only baseline.} We
state this deliberately, because the first comparison in isolation would
support the stronger and unwarranted claim that graph augmentation reduces
unsafe recommendations in general. On this benchmark it does not; it undoes
harm that structured record context causes.

This also sharpens our sparsity result. The graph
helps where the record actively misleads, which is a narrower and more
mechanistic account of when graph augmentation is worth its cost than
``knowledge helps when the record is thin''.

\subsection{Two limitations of the graph itself}

Both of the following are properties of the knowledge resource rather than the
routing system, and both surfaced only because decisions were adjudicated
individually.

First, our linearisation flattens conditional prohibitions into unconditional
ones. The graph stores qualifying notes---``contraindicated if
eGFR${<}30$''---but renders the fact to the model as ``Metformin is
contraindicated in Type 2 Diabetes Mellitus''. Read literally that is false;
metformin is first-line therapy for type 2 diabetes absent renal impairment. A
model that endorses the drug may therefore be clinically correct while
disagreeing with the fact it was given, and our measure will score it as a
violation. This inflates measured violations in an unknown but non-zero number
of cases and should be corrected in any deployment of this pipeline.

Second, the graph is incomplete in ways a clinician notices immediately. The
empagliflozin case above is a concrete instance: a drug with a well-known renal
risk profile carries no corresponding edge for a patient with impaired renal
function. Our hand-curated 25-disease graph encodes a fraction of what a
clinician would bring to the same decision.

Taken together these bound the interpretation. The measure reports
\emph{agreement with the graph's assertions}, which is what a faithfulness
check can establish; it is not a clinical safety audit, and should not be read
as one.

\subsection{Scope}

This covers one relation type of five, on a benchmark where violations are
rare---between eleven and eighteen per mode across 300 questions---so the
paired tests rest on six and seven discordant pairs respectively. Labels come
from a single annotator who is an author, blind to system identity and to
detector output but not to our hypotheses, and no inter-annotator agreement
study was performed. The direction and mechanism are consistent across the
three comparisons and with an independent observation recorded earlier in this
project's development, but the effect sizes are small-sample estimates.
